\originalchapter{恒星核反应、Gamow窗口与热反馈}
本章支撑8.901第9--10次课。太阳温度对应的热能远低于经典库仑势垒，经典高能尾不足以维持观测光度；量子隧穿、反应的弱作用瓶颈及核心的热反馈共同解释长期发光。需要Maxwell分布、相对坐标和WKB近似，具体实验$S$因子仍是输入。

\originalsection{反应率应使用相对速度}
两物种数密度$n_1,n_2$的反应数率为$n_1n_2\langle\sigma v_{\rm rel}\rangle$，同种粒子另除2以免双数。相对动能$E=\mu_r v_{\rm rel}^2/2$，$\mu_r$为约化质量，不能把恒星气体平均粒子质量$\mu$用在这里。两套Maxwell速度经过质心、相对速度变换，得到
\[
 \langle\sigma v_{\rm rel}\rangle
 =\sqrt{\frac8{\pi\mu_r}}(k_BT)^{-3/2}
 \int_0^\infty \sigma(E)\,E\,e^{-E/(k_BT)}\dd E .
\]
反应放热$Q$乘体积反应率、再除$\rho$给单位质量核产热；逃逸中微子携能必须从热化预算中扣除。PS4.8给相对Maxwell变量变换的完整证明，PS4.6用真实核质量核查不同通道的$Q$和正负号。

\originalsection{隧穿选择的是有限能窗}
非共振带电粒子截面写$\sigma=S(E)E^{-1}\exp[-\sqrt{E_G/E}]$。$E_G$由电荷乘积和约化质量决定，PS4.5用WKB积分推出其指数。核截面的缓变部分$S(E)$来自核物理测量，不能只由库仑势求定。

热尾与隧穿的联合指数$f=E/(k_BT)+\sqrt{E_G/E}$，极小条件
\[
 f'(E_0)=0,\quad E_0=[E_G(k_BT)^2/4]^{1/3},\quad
 f(E_0)=3E_0/(k_BT).
\]
二阶导数$f''(E_0)=3/(2E_0k_BT)$，所以主贡献的Gaussian近似为$\exp[-3(E-E_0)^2/(4E_0k_BT)]$。常用全宽$\Delta=4\sqrt{E_0k_BT/3}$是指数从峰值降$1/e$的左右跨度；它不是标准差或半高全宽。提高温度使窗口右移、降低指数抑制，较大电荷势垒使温度敏感性更强。遇到窄共振或快速变化$S$时不能继续当它常数作同一鞍点近似。

\originalsection{pp与CNO通道的净结果相同，时钟不同}
pp链从$p+p\to d+e^++\nu_e$开始，随后$d+p\to{}^3{\rm He}+\gamma$；两个$^3$He可生成$^4$He并放回两个质子，也可经更重中间核形成不同支链。第一步需要一个质子通过弱作用变为中子，缓慢的弱反应是长寿命的重要原因；仅发现质子能隧穿并不能决定整个链的速率。

CNO循环以已有碳氮氧为催化核，例如
\[
 {}^{12}{\rm C}(p,\gamma){}^{13}{\rm N}
 \xrightarrow{\beta^+}{}^{13}{\rm C}(p,\gamma){}^{14}{\rm N}
 (p,\gamma){}^{15}{\rm O}
 \xrightarrow{\beta^+}{}^{15}{\rm N}(p,\alpha){}^{12}{\rm C}.
\]
碳核回到起点，净同为四质子生成氦、正电子和中微子。净质量能约26.7 MeV包含正电子最终湮灭；实际留在星中的热能少一些，取决于分支中微子携能。CNO链的较高电荷势垒使它在较高温核心更重要，但还需初始催化核丰度；不是单凭恒星总质量就无条件决定唯一通道。

将实际速率在有限温区近似$\epsilon\propto T^\nu$可解释调节强弱，pp常较弱、CNO较强；$\nu$是局部对数斜率，不是全球恒定整数。核燃烧使温度升高，非简并核心膨胀并降温，通常构成负反馈。用给定温度幂律外推到任意温度、密度会丢失反应网络和组成变化。

\originalsection{三氦与重元素燃烧}
两个$\alpha$核先形成寿命很短的$^8$Be平衡中间态，再吸收$\alpha$进入碳的共振态；三体有效过程使单位质量产热大致含$\rho^2Y^3$而非二体反应的$\rho$。共振把指数能量位置改变，不能套用无共振两体$S$因子法给出同一个温度依赖。

原PS5.3所给三氦近似有$T_8^{-3}e^{-44/T_8}$，其局部斜率为$-3+44/T_8$，在$T_8=1$附近约41，极为敏感。简并电子压强对温度不敏感时，升温不能立即靠膨胀降温，产生氦闪；PS5.3实际积分还区分原题热容系数与单原子离子热容的差异。此恒密度局部闪模型不是完整恒星闪峰光变计算。

更大质量核心可依次经历碳、氖、氧、硅燃烧及复杂准平衡网络。铁峰核的结合能已接近最大，继续把它们聚变成更重核不再提供稳定的净核能来源。具体核素产额、爆发路径和实测速率没有被上述尺度律穷尽，后续核心坍缩章将用失去压强支持和能量预算说明终点。
